% !PRINTED 530
\weberpage{530}

Zerlegt man also das Integral (21) in seine beiden Bestandteile und setzt
\[
\int_0^\infty t^{2s-3}e^{-t}\,dt = \Gamma(2s-2) = \frac{\Gamma(2s-1)}{2(s-1)},
\]
so folgt:
\[
\mathop{\mathrm{Lim}}\limits_{s=1}\left(\int_0^\infty \frac{t^{2s-2}\,dt}{e^t-1} - \frac{\Gamma(2s-1)}{2(s-1)}\right) = -\Gamma'(1),
\]
und wenn man also die Entwickelung nach Potenzen von $s-1$:
\[
\Gamma(2s-1) = 1 + 2\Gamma'(1)(s-1) + \cdots
\]
einsetzt,
\weq{22}{\mathop{\mathrm{Lim}}\limits_{s=1}\left\{2\int_0^\infty \frac{t^{2s-2}\,dt}{e^t-1} - \frac{1}{s-1}\right\} = 0.}

Man erhält ferner durch Entwickelung nach Potenzen von $s-1$:
\weq{23}{\begin{aligned}
\frac{(2\pi)^{2s}A^{s-1}}{(4\pi\sqrt m)^{2s-1}\Gamma(s)^2}
&= \frac{\pi}{\sqrt m}\left[1+(s-1)\left(\log\frac{A}{4m}-2\Gamma'(1)\right)+\cdots\right]
\end{aligned}}
so daß nach (20):
\weq{24}{\begin{aligned}
P_s &= \frac{(2\pi)^{2s}A^{s-1}}{(4\pi\sqrt m)^{2s-1}\Gamma(s)^2}\left(2\int_0^\infty \frac{t^{2s-2}\,dt}{e_t-1} - \frac{1}{s-1}\right)\\ % [sic] printed $e_t-1$; equation (22) and the sense require $e^t-1$
&\quad + \frac{\pi}{\sqrt m(s-1)} + \frac{\pi}{\sqrt m}\left(\log\frac{A}{4m}-2\Gamma'(1)\right)+\cdots,
\end{aligned}}
worin die noch folgenden Glieder mit $s-1$ verschwinden. Daraus folgt:
\weq{25}{\mathrm{Lim}\left(P_s' - \frac{\pi}{\sqrt m(s-1)}\right) = \frac{\pi}{\sqrt m}\left(\log\frac{A}{4m}-2\Gamma'(1)\right).}
% [sic] printed $P_s'$; equation (24) defines $P_s$ without a prime.

Es bleibt noch der Bestandteil $Q_s$ zu untersuchen übrig, der für $s=1$ einen endlichen Grenzwert $Q$ erhält; diesen können wir ohne weiteres durch Einsetzen des Wertes $s=1$ bestimmen.

Es läßt sich nämlich in (10), solange $u$ und $y$ größer als Null sind, nach Einsetzen des Wertes $s=1$ die Integration
%CONT>>
